\textbf{Proof.}
Use the compact normalized set \(K_\kappa\) in (1) of the preceding proof.  Lemma lem:algebraic-box-feasibility-cad decides exactly whether each rational signed pivot-chart box meets \(K_\kappa\), including equality-only feasibility, and returns either a real-algebraic sample or a finite infeasibility certificate.

For \(e\in\{L,U\}\), define
\[
R_e=\Omega\Pi^{-1}Q(c_e)\Pi^{-1}\Omega,
\quad S(B)=B^{\top}\Omega B,
\quad T_e(B)=B^{\top}R_eB.
\]
The algebraic input yields rational bounds \(M_\Omega\ge\lVert\Omega\rVert_2\) and \(M_e\ge\lVert R_e\rVert_2\).  If \(d=\lVert B-C\rVert_F\), normalized representatives obey
\[
\lVert S(B)-S(C)\rVert_2\le2\sqrt{2n}M_\Omega d,
\qquad
\lVert T_e(B)-T_e(C)\rVert_2\le2\sqrt{2n}M_ed.
\tag{1}
\]
Because \(\lVert S(B)^{-1}\rVert_2\le(n\kappa)^{-1}\), the resolvent identity and (1) give an effectively computable rational \(\overline L\) satisfying
\[
|w(B)-w(C)|\le\overline L\lVert B-C\rVert_F
\quad(B,C\in K_\kappa).
\tag{2}
\]

Choose rational \(\rho>0\) with \(4\rho<\varepsilon\).  Uniformly bisect the finite chart cover until every retained box \(C\) has \(d_C\le\varepsilon\) and \(\overline Ld_C\le\rho\).  CAD deletes exactly infeasible boxes and supplies an algebraic point \(B_C\) in every feasible box.  Algebraic isolation gives rational bounds
\[
w(B_C)-\rho\le\underline q_C\le w(B_C)
\le\overline q_C\le w(B_C)+\rho.
\tag{3}
\]
Then \(\lambda_C=\underline q_C-\overline Ld_C\) is a valid box lower bound.  With \(L_\varepsilon=\min_C\lambda_C\), \(U_\varepsilon=\min_C\overline q_C\), and \(B_\varepsilon\) taken from an upper-bound minimizing box, a minimizer-containing box and (2)--(3) give
\[
v_{\mathrm{mm}}-2\rho\le L_\varepsilon\le v_{\mathrm{mm}}
\le w(B_\varepsilon)\le U_\varepsilon\le v_{\mathrm{mm}}+2\rho.
\tag{4}
\]
Thus the value gap is below \(\varepsilon\).  Pruning a feasible box only when its lower bound exceeds \(U_\varepsilon\) preserves every minimizer box, hence \(\mathcal A_{\mathrm{mm}}\subseteq\mathcal C_\varepsilon\); subdivision gives \(h_\varepsilon\le\varepsilon\).  CAD signs, algebraic samples, outward bounds, chart, norm, diameter, pruning, and mesh records form the asserted locally checkable finite certificate.  Weak polynomial inequalities retain the \(\kappa\)-boundary.  Equation (4) and the mesh bound prove convergence as \(\varepsilon\downarrow0\). \(\square\)